\section{The Chow ring}

A central object in intersection theory is the \emph{Chow ring} of an algebraic variety. The Chow ring encodes information about subvarieties of a given variety, and allows us to compute intersections of these subvarieties in a rigorous way. In this section, we give an introduction to the Chow ring, following \cite[Section 1.3]{FUL84} and \cite[Section 1.2]{EH16}.

Let $X$ be an algebraic scheme. A \emph{$k$-cycle} on $X$ is a finite formal sum 
\[
    \Sigma_i m_i [V_i]
\]
where the $V_i$ are $k$-dimensional subvarieties of $X$, and the $m_i$ are integers. The group of $k$-cycles on $X$, denoted $Z_k(X)$, is the free Abelian group generated by the $k$-cycles. To a $k$-dimensional subvariety $V$ of $X$ corresponds $[V]$ in $Z_k(X)$. The group of all cycles on $X$ is denoted 
\[Z(X) = \bigoplus_k Z_k(X),\]
and its elements are called \emph{cycles}.

\begin{remark}
    $(n-1)$-cycles are Weil divisors.
\end{remark}

For any $(k+1)$-dimensional subvariety $W$ of $X$, and any non-zero rational function $f \in K(W)^\times$, we can define a $k$-cycle $[\text{div}(f)]$ on $X$ by
\[
    [\text{div}(f)] = \sum_V \text{ord}_V(f) [V]
\]
where the sum is over all codimension-1 subvarieties $V$ of $W$. Here, $\text{ord}_V(f)$ is the order function on $K(W)^\times$ defined by the local ring $\mathcal{O}_{W,V}$. 

A $k$-cycle $\alpha$ is \emph{rationally equivalent to zero}, written $\alpha \sim 0$, if there are a finite number of $(k+1)$-dimensional subvarieties $W_i$ of $X$ and $f_i \in K(W_i)^\times$ such that
\[
    \alpha = \sum_{i} [\text{div}(f_i)].
\]
Since $[\text{div}(f^{-1})] = -[\text{div}(f)]$, the cycles rationally equivalent to zero form a subgroup of $\text{Rat}_k(X) \subseteq Z_k(X)$. The quotient group 
\[A_k(X) = Z_k(X)/\text{Rat}_k(X)\]
is called the \emph{$k$-th Chow group} of $X$, and its elements are called \emph{Chow classes}. We say that two $k$-cycles are \emph{rationally equivalent} if they lie in the same class in $A_k(X)$.

\begin{definition}
    The \emph{Chow group} of $X$ is the direct sum 
    \[A(X) = \bigoplus_k A_k(X).\]
\end{definition}

An element $\alpha \in A^k(X)$ is called a \emph{cycle class}. We say that a cycle is \emph{effective} if all its coefficients are non-negative. A cycle class is \emph{effective} if it can be represented by an effective cycle.

\begin{remark}
    In many cases, it will be natural to study codimension $r$ subvarieties of an $n$-dimensional variety $X$, and so we consider the $(n-r)$-th Chow group. We denote this as $A^r(X) = A_{n-r}(X)$. This group is called the \emph{codimension r Chow group} of $X$.
\end{remark}


% \todo[inline]{Include example 1.3.1 from Fulton's book?}

\begin{theorem}[\cite{EH16} Theorem-definition 1.5]\label{thm:chow}
    If $X$ is a smooth quasi-projective variety, then there is a unique product structure on $A(X)$ satisfying the condition:

    If two subvarieties $A$, $B$ of $X$ are generically transverse, then
    \[[A][B] = [A \cap B].\]
    This structure makes 
    \[A(X) = \bigoplus_{c=0}^{\dim X} A^c(X)\]
    into an associative, commutative ring, graded by codimension, called the \emph{Chow ring} of $X$.
\end{theorem}

% We say that two $k$-cycles are \emph{rationally equivalent} if their difference is a finite sum of divisors of rational functions on $(k+1)$-dimensional subvarieties of $X$. The group of $k$-cycles modulo rational equivalence is called the \emph{$k$-th Chow group} of $X$, and is denoted $A_k(X)$.

% For more details on the Chow ring, see \cite[Chapter 1]{FUL84}, or \cite[Section 1.2]{EH16}.



\subsection{Alternative definition of rational equivalence}
A more geometric definition of rational equivalence, and therefore perhaps more intuitive to some, is the following. 
\begin{definition}
    \cite{EH16} Let $\text{Rat}(X) \subseteq Z(X)$ be the subgroup generated by differences of the form
\[\langle \Phi \cap (\{t_0\} \times X) \rangle - \langle \Phi \cap (\{t_1\} \times X) \rangle\]
where $t_0,t_1 \in \mathbb{P}^1$ and $\Phi$ is a subvariety of $\mathbb{P}^1 \times X$ not contained in any fiber $\{t\} \times X$. We say that two cycles are \emph{rationally equivalent} if their difference lies in $\text{Rat}(X)$, and we say that two subschemes are rationally equivalent if their associated cycles are rationally equivalent.
\end{definition}
\cite[Section 1.6]{FUL84} shows that this definition is equivalent to the previous one.
\todo[inline]{Legg til bilde?}

    
\subsection{The moving lemma}
One of the fundamental results in intersection theory is the \emph{moving lemma}, which allows us to represent Chow classes by cycles that intersect transversely. In fact, this was historically the basis for a proof of \cref{thm:chow}. Generally, the moving lemma extends an idea from divisors, where one can find the intersection product of two divisors, by considering extended divisors using very ample divisors. 
\todo[inline]{Fiks?}
%For instance, to calculate the self-intersection of an exceptional divisor $E$ on a surface, one cannot find a transverse intersection of $E.E$, as $E$ cannot actually be moved, it is in fact the only representative of its line bundle. However, by considering a divisor $D = (E+A)-A$, one can find a transverse intersection of $E.D$.

\begin{theorem}[The moving lemma]\label{thm:moving}
    Let $X$ be a smooth quasi-projective variety.
    \begin{enumerate}[label=(\alph*)]
        \item For every $\alpha, \beta \in A(X)$ there are generically transverse cycles $A, B \in Z(X)$ with $[A] = \alpha$ and $[B] = \beta$.
        \item The class $[A \cap B]$ is independent of the choice of such cycles $A$ and $B$.
    \end{enumerate}
\end{theorem}
We refer to \cite[Theorem 1.6]{EH16} for a proof.
